\chapter{Research and Existing Solutions}
\label{ch:research}
% ============================================================

\section{Research Method / Approach}
\label{sec:research-method}
The literature was investigated in three directions corresponding to the three
bodies of work the project sits between: (i) federated learning for intrusion
detection, (ii) automotive in-vehicle network security, and (iii) TSN protocol
security, specifically gPTP. Sources were drawn from IEEE Xplore, ACM Digital
Library, ScienceDirect and arXiv, together with the relevant IEEE~802.1 standards
documents.

Each source was recorded against four fields (methodology, results, limitations,
and the design decision it influenced) so that the review produces design
consequences rather than a summary. Candidate sources were retained only where they
addressed at least one of: TSN protocol-native detection signals, federated
aggregation under non-IID clients, or automotive zonal architecture.

In addition, a comparison was performed against the departmental FYP repository
(Fall 2019 -- Fall 2025) to establish that the project does not duplicate prior
student work. The result of that comparison is reported in Section~\ref{sec:comparison}.

% [[FILL]] You stated your review is broader than the five papers in the deck.
% List the additional sources here with the same four fields, then extend Table 2.1.

\section{Existing Solutions / Products}

\begin{longtable}{p{0.16\textwidth} p{0.26\textwidth} p{0.24\textwidth} p{0.26\textwidth}}
\caption{Reviewed solutions, results and limitations}\\
\toprule
\textbf{Source} & \textbf{Methodology} & \textbf{Results} & \textbf{Limitations} \\
\midrule
\endfirsthead
\toprule
\textbf{Source} & \textbf{Methodology} & \textbf{Results} & \textbf{Limitations} \\
\midrule
\endhead
SAFID (2026) &
Selective federated learning, TDNet, domain adaptation, post-training quantization &
99.75\% accuracy on CICIDS2017; 95.75\% on Edge-IIoTset; $6.5\times$
inference-throughput improvement &
Evaluated on generic IDS datasets rather than automotive zonal TSN traces;
industrial segmented networks, not vehicle zonal \\
\addlinespace
Y\"ucel et al.\ (2026) &
Physical avionics TSN testbed; unsupervised anomaly detection using Local Outlier
Factor on gPTP timing features &
Detection at 500\,ns and 1000\,ns offsets; F1 remained above 0.95 even for 25\,ns and
50\,ns attacks &
Avionics testbed, not automotive zonal TSN; no federated learning; single
synchronisation domain; unsupervised framing driven by label scarcity that does not
apply to a simulation-generated dataset \\
\addlinespace
TSNZeek (2023) &
Zeek extension with Spicy packet parsing and rule-based detection for TSN protocols &
$\approx$5\% CPU overhead; 25\% CPU at 50\,Mbit/s, 45\% at 250\,Mbit/s; up to 2\,ms
processing lag; detected all seven tested attacks &
Does not address gPTP timing attacks, federated learning, or automotive zonal
deployment; rules are stateful and centralized \\
\addlinespace
Sorescu et al.\ (2026) &
Model-based TARA using STRIDE and ISO/SAE 21434 assumptions &
151 distributed, 259 domain-based and 270 zonal threat rows &
Threat counts are model outputs, not measured attack probabilities or IDS results \\
\addlinespace
Adil et al.\ (2026) &
Federated collaborative learning, NTTE traffic encoding, TCNnet, INT8 quantization,
TIMM safety modelling &
Up to 99.98\% accuracy; model size reduced from 2.81\,MB to 0.23\,MB; collision
reduction up to 97.9\% &
Real-time vehicle integration, rare attacks and protocol-specific anomalies remain
unresolved \\
\bottomrule
\end{longtable}

\section{Comparison of Existing Solutions}
\label{sec:comparison}

Three dimensions are jointly required by this project: federated training, an
automotive zonal setting, and detection built on TSN protocol-native signals.
Table~\ref{tab:dims} evaluates the reviewed literature and the closest prior
departmental FYPs against those dimensions.

\begin{longtable}{p{0.34\textwidth} c c c}
\caption{Coverage of the three required dimensions}\label{tab:dims}\\
\toprule
\textbf{Work} & \textbf{FL} & \textbf{Automotive / zonal} & \textbf{TSN-native signals} \\
\midrule
\endfirsthead
\toprule
\textbf{Work} & \textbf{FL} & \textbf{Automotive / zonal} & \textbf{TSN-native signals} \\
\midrule
\endhead
SAFID (2026)                & Yes & No & No \\
Y\"ucel et al.\ (2026)      & No  & No (avionics) & Yes \\
TSNZeek (2023)              & No  & No (industrial) & Partial (rule-based) \\
Adil et al.\ (2026)         & Yes & Yes (CAN/embedded) & No \\
F25-151-R-DIDS-IoV (FYP)    & Yes & Yes (IoV, not zonal) & No \\
F25-209-R-BastionFed (FYP)  & Yes & No & No \\
F25-378-D-AlertEye (FYP)    & Yes & No & No \\
F25-178-R-SafeCharge (FYP)  & No  & EV charging, off-vehicle & No \\
F25-168-D-NeuroDrive (FYP)  & No  & Yes (OBD-2 bus) & No \\
\midrule
\textbf{CARAZONE (this work)} & \textbf{Yes} & \textbf{Yes} & \textbf{Yes} \\
\bottomrule
\end{longtable}

No reviewed work covers all three dimensions. One prior departmental project
(DIDS-IoV) covers two. Across the departmental repository from Fall~2019 to
Fall~2025, no project addresses TSN, gPTP, FRER, 802.1 deterministic networking or
automotive Ethernet; the single automotive in-vehicle network security project
targets the CAN bus.

\section{Limitations and Gaps}
\label{sec:limitations-gaps}

The reviewed work exhibits four limitations relevant to this project:

\begin{enumerate}
  \item \textbf{Feature space.} Enterprise, IoT and industrial IDS work operates on
        protocol-agnostic statistical traffic features, because those networks offer
        no configured schedule or synchronisation structure to compare against. In a
        deterministic network, the expected value of a timing observation is known in
        advance from configuration. No reviewed work exploits this asymmetry in an
        automotive zonal setting.
  \item \textbf{Domain transfer.} Work that does use gPTP timing signals is situated
        in avionics, with a single synchronisation domain and a different traffic
        profile from a vehicle zonal backbone.
  \item \textbf{Federated evaluation realism.} Federated IDS work commonly evaluates
        on artificial IID partitions. Zone-level heterogeneity in a vehicle is
        structural, since a perception zone and a body zone differ in traffic density,
        stream composition and criticality, and is not represented by a random
        split.
  \item \textbf{Operational framing.} Several works report classification quality
        without detection latency, alert volume at realistic base rates, or the
        effect of monitoring on the protected traffic it observes.
\end{enumerate}

\textbf{Gap statement.} No existing work provides federated, zone-level intrusion
detection for automotive TSN backbones using protocol-native timing signals, with
evaluation that reports detection quality, attribution, latency and federated
behaviour under structurally heterogeneous clients.

\section{Proposed Solution}
CARAZONE addresses the gap with a centralized cross-silo horizontal federated
learning architecture in which each zone controller is a client. Each client
observes local gPTP state passively, computes windowed protocol-native features,
and trains a shared compact model architecture. The central vehicle compute unit
aggregates parameters using FedAvg and redistributes the global model. Clients then
perform inference locally and emit attributed alerts.

Two design choices follow directly from the gap:

\textbf{Protocol-native features.} Features are computed as residuals against values
that TSN configuration makes known in advance: deviation of a node's offset from
its siblings in the synchronisation tree, deviation of sync message arrival from the
configured interval, and the variance of measured link delay. These are not available
to a generic traffic-feature detector, and the controlled comparison against such a
detector is the project's primary experiment.

\textbf{Federation across zones rather than vehicles.} Clients are zones within one
vehicle, so the aggregation server is the vehicle's own central compute unit. The
motivation is not inter-organisational privacy but bandwidth: streaming full timing
traces from every zone to a central analyser is unbounded best-effort traffic
competing with the scheduled streams the system protects, whereas periodic exchange
of fixed-size model parameters is a bounded, schedulable traffic pattern.

\section{Complex Computing Problem}

The project addresses a complex computing problem under the Seoul Accord
characterisation. Complexity does not arise from combining TSN, machine learning and
federated learning; it arises from the interaction between timing guarantees,
low-observable attacks, heterogeneous distributed training and operational
constraints.

\begin{longtable}{p{0.3\textwidth} p{0.64\textwidth}}
\toprule
\textbf{Characteristic} & \textbf{Manifestation in this project} \\
\midrule
\endhead
Conflicting requirements &
Detection sensitivity competes with false-alarm volume, inference cost on a
resource-constrained zone controller, federated communication overhead, and the
requirement that monitoring must not perturb the protected traffic it observes. \\
\addlinespace
No obvious solution &
A stealthy delay attack can be constructed to remain inside the envelope of benign
oscillator drift, so no fixed threshold separates the two; and federated aggregation
may degrade rather than improve performance under genuinely non-IID clients. \\
\addlinespace
Depth of knowledge &
Requires IEEE~802.1AS synchronisation mechanics, time-aware scheduling, discrete-event
network simulation, timing-feature design, and federated optimisation. \\
\addlinespace
Beyond standard practice &
Neither IEEE~802.1AS nor IEEE~802.1Qbv defines a security or detection layer. A
routine IDS or a routine FL implementation provides neither timing-aware attribution
nor evaluation of monitoring interference. \\
\addlinespace
Significant consequences &
Missed detection permits forced degradation of safety-relevant functions; excessive
false alarms make the system operationally unusable at fleet scale. \\
\addlinespace
Interdependence &
Synchronisation behaviour, feature extraction, local training, aggregation strategy
and deployment inference cost are mutually constraining; a change in window size
propagates to detection latency, model size and communication volume. \\
\bottomrule
\end{longtable}

\section{Technical Challenges}
\label{sec:technical-challenges}

\begin{enumerate}
  \item \textbf{Observability of internal protocol state.} Several informative
        features derive from the synchronisation servo's internal state rather than
        from any field visible in a packet capture. Whether the simulation framework
        exposes these, or whether explicit instrumentation is required, must be
        verified per feature before the feature is committed.
  \item \textbf{Benign/attack separability.} Benign oscillator behaviour includes
        thermal and ageing drift whose magnitude can exceed that of a deliberately
        stealthy attack. Detection must therefore rely on structural rather than
        magnitude-based discrimination.
  \item \textbf{Grandmaster-level blind spot.} Detection based on comparison between
        siblings in the synchronisation tree is, by construction, insensitive to a
        perturbation that moves the entire tree together. This was confirmed
        empirically (Section~\ref{sec:edge-cases}) and is recorded as a scope boundary: the system
        addresses link-local timing integrity, not time-source integrity.
  \item \textbf{Structural client heterogeneity.} For the federated comparison to be
        meaningful, zones must differ in traffic profile. A symmetric topology makes
        federated, centralized and local-only training converge to equivalent
        models and renders the comparison null by construction.
  \item \textbf{Class imbalance and evaluation validity.} Attack windows are a small
        minority of all windows. Accuracy is therefore uninformative, and train/test
        splitting must occur at scenario and seed granularity rather than at window
        granularity to prevent leakage between adjacent correlated windows.
\end{enumerate}

\section{Research Findings and Design Decisions}
\label{sec:research-findings}

\begin{longtable}{p{0.24\textwidth} p{0.34\textwidth} p{0.36\textwidth}}
\toprule
\textbf{Finding} & \textbf{Source} & \textbf{Design decision} \\
\midrule
\endhead
Automotive TSN can be represented with domains, gateways, periodic streams and
scheduling in a discrete-event simulator &
Zhou et al.\ (2021) &
Use OMNeT++/INET rather than building a simulator from scratch; keep the topology
small and seeded \\
\addlinespace
Clock offset, port state and synchronisation message behaviour expose automotive
Ethernet synchronisation attacks &
Fotouhi et al.\ (2023) &
Commit offset, offset trend and synchronisation-interval features as the core
feature family \\
\addlinespace
Some PTP attacks are stealthy and consume very little bandwidth &
Berardi et al.\ (2023) &
Do not rely on packet-volume features; make the generic-traffic-feature detector the
control arm rather than the baseline to beat on volume \\
\addlinespace
gPTP timing features support detection at small offset magnitudes, with F1 above 0.95
at 25\,ns and 50\,ns &
Y\"ucel et al.\ (2026) &
Adopt comparable timing features; adopt LOF as an additional baseline so results are
comparable to published work \\
\addlinespace
Unsupervised framing in prior work is a response to label scarcity in real
deployments &
Y\"ucel et al.\ (2026) &
Use supervised multiclass classification, since simulation supplies ground-truth
labels; this also yields a model with trainable weights, which plain LOF does not,
and is therefore FedAvg-compatible without a non-standard aggregation scheme \\
\addlinespace
Centralized, local-only and FedAvg comparison exposes the trade-off between federated
performance and communication cost &
Digregorio et al.\ (2024) &
Commit all three arms as the federated evaluation; retain FedProx as an extension \\
\addlinespace
Random IID partitions and accuracy-only evaluation overstate federated performance &
Shahsavari et al.\ (2026) &
Use structurally heterogeneous zone profiles, split by scenario and seed, and report
worst-client performance rather than only the mean \\
\addlinespace
Loss or manipulation of synchronisation causes clock divergence with a characteristic
recovery profile &
Kakade et al.\ (2022) &
Use divergence and post-attack recovery behaviour as attack-validation signals before
any model is trained \\
\bottomrule
\end{longtable}

% ============================================================